% =============================================================================
%  Section 7 — Responsible Disclosure and Ethics Statement
% =============================================================================
\section{Responsible Disclosure and Ethics Statement}
\label{sec:ethics}

\paragraph{System provenance and authorship.}
Neuron is a research prototype developed by the authors of this paper; it
is not a third-party commercial product, and the authors are its sole
maintainers. The system name ``Neuron''\footnote{The name ``Neuron'' is
chosen as a research-system identifier. An unrelated open-source PHP web
application framework sharing a similar name exists; it has a completely
different architecture (PHP web framework vs.\ the Java-based governance
runtime described here), and the two projects share no code, design, or
maintainers.} is used solely for anonymity during double-blind review and
does not correspond to any deployed commercial product.

\paragraph{Vulnerability disclosure.}
The two vulnerability classes identified in this paper---coupled
resource--security failure (Section~\ref{sec:bg-lim}) and the LIM-specific
spawn-time attack surface (Section~\ref{sec:bg-attack})---were discovered
by the authors during the design and implementation of Neuron. Because
Neuron is the authors' own system, the ``disclosure'' is self-disclosure
of design flaws that were identified and remediated during development,
not a third-party vulnerability report. No CVE was filed, because the
vulnerabilities are design-level properties of a research prototype, not
implementation bugs in a deployed product.

\paragraph{Fix timeline.}
Both vulnerability classes were remediated in the Neuron codebase prior to
paper submission:
\begin{itemize}
    \item \textbf{Coupled resource--security failure}: The unified audit
    chain (\texttt{UnifiedAuditLog}), per-turn trace identifier
    (\texttt{TraceContext}), and \IO{} circuit-breakers
    (\texttt{VfsRateLimiter}, \texttt{EventBusRateLimiter}) were
    implemented in the initial runtime build. The coupled enforcement loop
    (Section~\ref{sec:arch-coupled}) has been present since the first
    evaluated version.
    \item \textbf{Spawn-time privilege escalation}: The
    \texttt{SpawnPrivilegeContext} (privilege non-increase) and
    tenant-identity propagation via \texttt{CallerContext} were
    implemented before the evaluation of Section~\ref{sec:eval-cross}. The
    depth-aware \texttt{EscalationPolicy} was added in the same commit.
    No production deployment existed prior to the fix.
\end{itemize}

\paragraph{Ethics and broader impact.}
This work studies security mechanisms for multi-tenant LLM agent runtimes.
The governance scenarios are evaluated in a controlled research environment
against the authors' own system; no third-party systems, users, or data
were involved. The evaluation harness and probe payloads are released as
part of the replication package to enable independent verification and
future defense research; they do not target any specific deployed system
or individual. The supplementary real-LLM experiment
(Section~\ref{sec:eval-supplement}) uses commercial LLM APIs under their
respective terms of service, with API costs borne by the authors; no
personally identifiable information was processed, and the adversarial
prompts were confined to the controlled experimental context.

\paragraph{Potential for harm and mitigations.}
The attack techniques described in this paper (resource-contention
flooding, spawn-time privilege escalation, natural-language escalation)
are applicable to any multi-agent framework that lacks spawn-time
privilege propagation. To mitigate misuse, we delay the release of
fully weaponized attack scripts until after publication, and the released
replication package focuses on the \emph{defense} mechanisms and
\emph{evaluation harness} rather than on ready-to-run exploit code. The
naive baseline experiment (Section~\ref{sec:eval-naive}) demonstrates
that forward-path defenses are insufficient, which we hope motivates the
community to adopt structural enforcement rather than to exploit the gap.

\paragraph{Data availability.}
The evaluation data, experiment scripts, and raw LLM responses are
preserved in the replication package: statistical analysis
(\texttt{statistical\_analysis.json}), annotation agreement
(\texttt{annotation\_agreement.csv}), naive baseline
(\texttt{naive\_baseline.csv}), and second-system reproduction
(\texttt{second\_system\_reproduction.csv}). To support reproducibility
during double-blind review, an anonymized artifact bundle containing
the raw result files (CSV/JSON), the evaluation harness, the
capability-surface classifier, and the TLA$^+$ specification is
available via the anonymous supplementary archive referenced in the
cover letter; the full system source and the weaponized attack scripts
will be released post-publication under an open-source license. The
artifact bundle allows reviewers to re-run the deterministic
experiments (Scenario~6, capability classifier, joint-decision,
starvation) and to inspect the raw LLM responses underlying the
stochastic experiments, though re-running the LLM trials requires the
reviewer's own API credentials.
